\documentclass[aspectratio=169]{beamer}
\usepackage[utf8]{inputenc}
\usepackage[polish]{babel}
\usepackage{amsmath,amssymb}
\usepackage{graphicx}

\usetheme{Madrid}
\usecolortheme{whale}

\title[Fizyka Współczesna i Klasyczna]{Wybrane Zagadnienia z Fizyki}
\subtitle{Podstawowe prawa, mechanika kwantowa i teoria względności}
\author{Projekt Edukacyjny}
\date{\today}

\usepackage[utf8]{inputenc}
\usepackage[T1]{fontenc}
\begin{document}

% Slajd 1: Tytułowy
\begin{frame}
    \titlepage
\end{frame}

% Slajd 2: Spis treści
\begin{frame}{Spis Treści}
    \tableofcontents
\end{frame}

\section{Wprowadzenie}

% Slajd 3: Listy punktowane i numerowane
\begin{frame}{Czym zajmuje się fizyka?}
    Fizyka to nauka przyrodnicza badająca fundamentalne i uniwersalne prawa przyrody.
    
    \vspace{0.3cm}
    Główne działy klasyczne:
    \begin{itemize}
        \item Mechanika klasyczna i grawitacja
        \item Termodynamika i fizyka statystyczna
        \item Elektrodynamika i optyka
    \end{itemize}
    
    \vspace{0.3cm}
    Etopy rozwoju poznania naukowego:
    \begin{enumerate}
        \item Obserwacja zjawiska i zbieranie danych
        \item Sformułowanie hipotezy badawczej
        \item Przeprowadzenie eksperymentu weryfikacyjnego
        \item Utworzenie teorii matematycznej
    \end{enumerate}
\end{frame}

\section{Mechanika Klasyczna}

% Slajd 4: Wzory matematyczne
\begin{frame}{Mechanika Klasyczna i Dynamika}
    Podstawą mechaniki klasycznej są zasady dynamiki Newtona oraz zasady zachowania.
    
    \begin{block}{Druga zasada dynamiki Newtona}
        Jeżeli na ciało działa niezrównoważona siła, to porusza się ono z przyspieszeniem wprost proporcjonalnym do działającej siły i odwrotnie proporcjonalnym do masy ciała:
        \begin{equation}
            \vec{F} = m \cdot \vec{a}
        \end{equation}
    \end{block}
    
    \begin{exampleblock}{Zasada zachowania pędu}
        W układzie izolowanym całkowity pęd układu ciał pozostaje stały w czasie:
        \begin{equation}
            \sum \vec{p}_i = \text{const}
        \end{equation}
    \end{exampleblock}
\end{frame}

\section{Termodynamika}

% Slajd 5: Bloki informacyjne
\begin{frame}{Termodynamika i Ciepło}
    Termodynamika opisuje procesy związane z ciepłem, temperaturą oraz pracą.
    
    \begin{block}{Pierwsza zasada termodynamiki}
        Zmiana energii wewnętrznej układu $(\Delta U)$ jest równa sumie dostarczonego ciepła $(Q)$ oraz wykonanej pracy $(W)$ nad układem:
        \begin{equation}
            \Delta U = Q + W
        \end{equation}
    \end{block}

    \begin{alertblock}{Ostrzeżenie / Druga zasada}
        Ciepło nie może samorzutnie przepływać z ciała o niższej temperaturze do ciała o wyższej temperaturze bez wkładu pracy zewnętrznej.
    \end{alertblock}
\end{frame}

\section{Elektrodynamika}

% Slajd 6: Układ dwóch kolumn
\begin{frame}{Elektrodynamika i Pola}
    Elektrodynamika bada wzajemne oddziaływania ładunków elektrycznych oraz pól elektromagnetycznych.
    
    \begin{columns}[T]
        \column{0.48\textwidth}
        \textbf{Pole Elektrostatyczne}
        \begin{itemize}
            \item Prawo Coulomba opisuje siłę między ładunkami.
            \item Potencjał elektryczny i natężenie pola.
            \item Kondensatory jako magazyny energii.
        \end{itemize}
        \begin{equation}
            F = k \cdot \frac{q_1 q_2}{r^2}
        \end{equation}
        
        \column{0.48\textwidth}
        \textbf{Pole Magnetyczne}
        \begin{itemize}
            \item Prąd elektryczny wytwarza pole magnetyczne.
            \item Siła Lorentza działająca na ładunek.
            \item Indukcja elektromagnetyczna Faradays.
        \end{itemize}
        \begin{equation}
            \vec{F}_L = q(\vec{E} + \vec{v} \times \vec{B})
        \end{equation}
    \end{columns}
\end{frame}

\section{Fizyka Kwantowa}

% Slajd 7: Fizyka kwantowa
\begin{frame}{Wprowadzenie do Fizyki Kwantowej}
    Na przełomie XIX i XX wieku klasyczna fizyka przestała tłumaczyć zjawiska w skali atomowej.
    
    \begin{block}{Hipoteza Plancka i dualizm korpuskularno-falowy}
        Energia promieniowania elektromagnetycznego jest emitowana w postaci porcji (kwantów) zwanych fotonami:
        \begin{equation}
            E = h \cdot \nu
        \end{equation}
        gdzie $h$ to stała Plancka, a $\nu$ to częstotliwość fali.
    \end{block}
\end{frame}

% Slajd 8: Równanie Schrödingera
\begin{frame}{Równanie Falowe Schrödingera}
    Fundamentalne równanie mechaniki kwantowej opisujące ewolucję stanu kwantowego w czasie:
    
    \begin{block}{Równanie Schrödingera (zależne od czasu)}
        \begin{equation}
            i \hbar \frac{\partial}{\partial t} \psi(\mathbf{r}, t) = \hat{H} \psi(\mathbf{r}, t)
        \end{equation}
    \end{block}
    
    \begin{itemize}
        \item $\psi$ – funkcja falowa cząstki
        \item $\hbar$ – zredukowana stała Plancka ($\hbar = \frac{h}{2\pi}$)
        \item $\hat{H}$ – operator Hamiltona (energia całkowita)
    \end{itemize}
\end{frame}

\section{Teoria Względności}

% Slajd 9: Szczególna teoria względności
\begin{frame}{Szczególna Teoria Względności}
    Stworzona przez Alberta Einsteina w 1905 roku, opiera się na dwóch postulatach:
    
    \begin{enumerate}
        \item Zasada względności – prawa fizyki są takie same we wszystkich układach inercjalnych.
        \item Stałość prędkości światła w próżni ($c$) dla każdego obserwatora.
    \end{enumerate}
    
    \begin{exampleblock}{Najsłynniejszy wzór fizyczny}
        Równoważność masy i energii:
        \begin{equation}
            E = m c^2
        \end{equation}
    \end{exampleblock}
\end{frame}

% Slajd 10: Podsumowanie
\begin{frame}{Podsumowanie}
    Fizyka stanowi fundament rozwoju nowoczesnej technologii i naszego zrozumienia wszechświata.
    
    \vspace{0.4cm}
    \begin{center}
        \textbf{Dziękuję za uwagę!} \\
        \vspace{0.2cm}
        Pytania i dyskusja.
    \end{center}
\end{frame}

\end{document}
